%Conclusiones 
Con el cumplimiento del Módulo de Ruteo Seguro para Alertify se evidencia el que a través de los objetivos planteados se consolidó una solución tecnológica frente a la problemática de la inseguridad urbana. Desde la obtención de requisitos, la adopción del marco de trabajo y la formalización mediante Historias de Usuario permitió una alineación precisa ante las necesidades de prevención del ciudadano y las especificaciones técnicas del sistema. Este enfoque iterativo garantizó que el diseño arquitectónico y funcional se mantuviera enfocado en la integridad física del peatón, validando la eficacia de las metodologías de desarrollo iterativas para resolver problemas de impacto social.

En el ámbito arquitectónico y de procesamiento de datos, la estructuración del esquema de base de datos relacional en SQL Server demostró ser un pilar fundamental para la viabilidad del proyecto. La utilización de tipos de datos espaciales nativos permitió modelar la red vial de la ciudad como un grafo complejo, donde la variable de peligrosidad se integró matemáticamente como un peso restrictivo en las aristas. La lógica de negocio, centralizada en la API desarrollada con el entorno de ejecución Nest.js, procesó estos grafos con precisión algorítmica para generar trayectorias óptimas de bajo riesgo. Si bien la topología de despliegue híbrida temporal con el motor de cálculo operando en un entorno local y la persistencia de datos alojada remotamente en la infraestructura de Microsoft Azure certificó la interoperabilidad de los módulos, también expuso las limitaciones inherentes a la latencia de red en arquitecturas distribuidas de esta naturaleza.

La calidad  y la viabilidad del software fueron ratificadas a través de una rigurosa fase de validación técnica y humana. La construcción del cliente nativo en Kotlin proporcionó una interfaz móvil fluida y altamente interactiva que logró abstraer la complejidad matemática del sistema. Con las pruebas se certificó la evasión proactiva de zonas de riesgo y la capacidad del sistema para ejecutar recálculos dinámicos ante desvíos imprevistos. Por su parte, la evaluación empírica de usabilidad permite concluir que, a pesar del cuello de botella en el procesamiento derivado de las peticiones remotas a la base de datos (con tiempos de espera oscilantes entre los 50 y 60 segundos), los usuarios perciben la herramienta como intuitiva, confiable y de alto valor estratégico para la toma de decisiones durante su movilidad urbana.